
\section{Is It Progress?}

Editorial in "Folkebladet," August 11, 1897.

It is not enough that there is a difference between the Free Church and the church bodies, or between congregational freedom and church government. There is a question that is still more important, and it is this: Is the Free Church now also truly a progress?

For those who believe that the congregation is about the same thing as rebellion and disorder and anarchy, while it is the church body with its wise government that brings a little order into the confusion, it is self-evident that the free-church tendency is a nuisance and a relapse into chaos. They reason in this way: in the church body the congregations have with great labor been gathered together, just as a builder laboriously drags the stones together for a building. And just as the one who tears the stones apart again and scatters them over the field is a vandal, so much more is he a vandal of the worst kind who works for the congregations to be independent and self-governed.*

These ways of speaking do indeed sound so reasonable and sensible that one might think the matter were thereby settled. One does not think of the fact that, if that sort of language could settle the matter, then it was also settled that the old Roman Church, with the Pope at its head, had the noblest and best right of all church bodies, and that it could with good reason complain that such rebels as Luther and others had disturbed the glorious building that had been raised with such great diligence and such great labor.

The question is this: Does it go forward when it goes from the state church toward the free congregation, so that it is progress so long as it goes that way? Or is it progress when the development goes from the free congregation toward an ever stronger church government, until at last this is everything and the congregation nothing, or at least nothing other than a mere supporting foundation?

One must not mix these two opposite things together. Either the free congregation is the goal,* and then one must work for it; or else the strong church body is the goal,* and then the task is to work for that, however it may go with the congregation and its freedom.

What shall we choose? History has indeed pronounced its judgment, so that we really need not doubt. The Catholic Church advanced great and strong in the delusion that the true progress was to build the Church ever higher, with the congregation deep below as a good foundation, until the Church became "a tower with its top in heaven" that gathered all Christians together and preserved them from dispersion. But God struck that proud church-building a crushing blow, and the Christians were scattered abroad as they were set free from the crushing absolute rule. Was it progress or decline when Luther shook the mighty church walls until the cracks appeared everywhere?

It was progress. And though the progress was lamed in its course when secular princes inherited the Pope’s power, yet it is again progress when the free people, as here among us, forbids the civil authority to govern the people’s religion and rule over consciences by worldly power and law.

But where the state church is no more, there the so-called "church bodies" now try to fill its place and preserve the old church order, which is the congregation’s tutelage.* An annual meeting is arranged, which grants the funds, and there is a chairman who appoints the ministers; and the congregations get less and less influence upon their own affairs and those of the church body, the greater the body becomes.

It is progress when congregational freedom is actually carried through, so that the congregations gain the right to determine their own expenses, choose their own ministers, and keep them undisturbed by the "church body" and its chairman. If the "church body" could have remained faithful to its task and defended the congregations’ freedom, then it would surely have been a long time before any doubt had been raised about its legitimacy.* But when the "church body" set power in the place of right and no longer took any account of the congregations’ freedom, but drove ministers from their congregations, cast congregational delegates out of the annual meeting, and expelled congregations from the body, then it became clear that here a tyrannical power had been established, under which the congregations could bow only with shame and harm.

The expelled congregations, and many with them, could not cease to exist; neither could they begin a new "body," when they had precisely been expelled by an old one. They had no other way out than to go straight forward to the full realization of congregational freedom.* In that there is a progress, and one that is all the more secure because it is biblical and Lutheran. 

Georg Sverdrup

NOTES ON UNCERTAIN READINGS

* "Anarki[?]sme" has been understood as "Anarkisme" = "anarchy."
* "frikirkelige Ste[v?]" has been rendered contextually as "the free-church tendency." The damaged reading appears to indicate a word such as "stræv" or "strev," but the sense is clearly a free-church movement or tendency.
* "æd[l?]ste" has been understood as "ædleste" = "noblest."
* Asterisks in the source have been retained where they appear to mark emphasis or editorial significance.
* "Stattefundament" has been translated contextually as "a mere supporting foundation," since the exact printed form appears unusual or damaged, but the sense is clear: the congregation reduced to a passive supporting base.
* "Kunde 'Kirkesamfundet' blevet sin Opgave tro" has been understood as "If the 'church body' could have remained faithful to its task." The source likely reflects a damaged or archaic printed form where the intended sense is evident from context.
* "Menighedsfrihedens fuldstændige Gjennemførelse" has been rendered as "the full realization of congregational freedom" to preserve both completion and principled implementation.
